\documentclass[11pt]{article}

\usepackage[a4paper,margin=29mm]{geometry}
\usepackage[T1]{fontenc}
\usepackage{lmodern}
\usepackage{microtype}
\usepackage{amsmath,amssymb,amsthm,mathtools}
\usepackage{booktabs}
\usepackage{array}
\usepackage{float}
\usepackage{enumitem}
\usepackage[hidelinks]{hyperref}

\newtheorem{theorem}{Theorem}[section]
\newtheorem{proposition}[theorem]{Proposition}
\newtheorem{corollary}[theorem]{Corollary}
\newtheorem{lemma}[theorem]{Lemma}
\theoremstyle{definition}
\newtheorem{definition}[theorem]{Definition}
\newtheorem{example}[theorem]{Example}
\theoremstyle{remark}
\newtheorem{remark}[theorem]{Remark}

\newcommand{\PG}{\mathrm{PG}}
\newcommand{\PGL}{\mathrm{PGL}}
\newcommand{\PSL}{\mathrm{PSL}}
\newcommand{\F}{\mathbb{F}}
\newcommand{\PP}{\mathbb{P}}
\newcommand{\Z}{\mathbb{Z}}
\newcommand{\Grundy}{\mathcal{G}}
\newcommand{\Cyc}{C}
\newcommand{\Mob}{M}
\DeclareMathOperator{\Aut}{Aut}
\DeclareMathOperator{\Fix}{Fix}
\DeclareMathOperator{\Sch}{Sch}
\DeclareMathOperator{\Cay}{Cay}
\DeclareMathOperator{\mex}{mex}
\DeclareMathOperator{\ord}{ord}
\DeclareMathOperator{\wt}{wt}

% Phase notes for the C264 execution chain (C306--C311).
% C309 must set \draftnotesfalse; no note may survive into the release build.
\newif\ifdraftnotes
\draftnotestrue
\newcommand{\phasenote}[2]{%
  \ifdraftnotes
    \par\smallskip\noindent
    \fbox{\parbox{0.96\linewidth}{\footnotesize\textbf{[#1 owed]}\ #2}}%
    \par\smallskip
  \fi}

\title{Node Kayles on Conic Schreier Graphs:\\
Dihedral and Polyhedral Templates}
\author{Tavis Rudd}
\date{Working draft --- July 2026}

\begin{document}
\maketitle

\begin{abstract}
Let $C$ be a nonsingular conic in $\PG(2,q)$ with $q$ odd. An off-conic point $x$ induces an
involution $\sigma_x$ of $C\cong\PP^1(q)$ by exchanging the two points in which each secant
through $x$ meets $C$. In a normal-play capacity-two avoidance game, selected off-conic points
forbid the conic points on their pairwise joining lines; if subsequent play is restricted to the
conic, the legal continuations form a Node-Kayles position. We prove one reduction and give three
applications of it. The reduction holds for any finite involution-generated action: the residual
graph is the fixed-point-deleted Schreier graph, and its Grundy value is the xor of the
transitive-template nimbers weighted only by orbit-multiplicity parity. The applications classify
every tame subgroup a legal selected set can induce. Two selected points realize every dihedral
group $D_{2m}$, in either parity of the rotation order, with cycle and path templates and hence
Dawson's-chess nimbers. Three selected points force the central involution and the even-rotation
groups $D_{4n}$, with empty, $K_4$, Möbius-ladder, prism, and ordinary-ladder templates of value in
$\{0,1\}$. The remaining tame possibilities are the finite polyhedral groups $S_4$ and $A_5$. The
resulting Grundy classification is an explicit congruence function of $q$, periodic in $q$ within
each torus family, and admits a converse realization theorem. Among admissible prime fields the
relative density of P-positions is one half for the even families, and one for every odd-order
dihedral group, where the game is always a P-position.
\end{abstract}

\noindent\textbf{Keywords.} Node Kayles; Sprague-Grundy theory; Schreier graph; conic; projective
line; dihedral group; $\PGL_2(q)$; Dawson's chess.

\smallskip
\noindent\textbf{2020 Mathematics Subject Classification.} 91A46, 05C57, 51E20, 20B25.

\phasenote{C308}{Keyword list, MSC codes, and the novelty calibration of the abstract are
C308 deliverables; the entries above are placeholders inherited from the source manuscript's
subject matter.}

\section{Introduction and main theorem}
\label{sec:intro}

Node Kayles is the impartial game in which a move at a graph vertex deletes its closed
neighbourhood. This paper studies a structured family of Node-Kayles positions arising from a
finite-geometric capacity-two game on a conic, and classifies their Grundy values whenever the
configuration of selected points induces a tame dihedral or polyhedral subgroup of $\PGL_2(q)$.

The general reduction of late-stage Nofil positions to Node Kayles is known
\cite{HugganHuntemannStevens2021}. The correspondence between off-conic points and involutions in
$\PGL_2(q)$, and the geometry of triples of such involutions, are also classical; Tranchida
\cite{Tranchida2025} recently developed that correspondence in detail for incidence geometries and
hypertopes. Our contribution starts after that correspondence. Saturation is detected by fixed
points of pair products, so the conic-only residual is a \emph{fixed-point-deleted Schreier graph};
its orbit templates and Grundy values can then be classified completely in the tame case.
Tranchida does not consider this residual graph or its Node-Kayles value.

\subsection{The reduction and its three applications}

The paper has one structural engine and three applications of it. The engine
(Theorem~\ref{thm:orbit-template}) says that for any finite group $G$ acting on a finite set
$\Omega$ and generated by a finite set $T$ of involutions, the residual Grundy value depends on
$\Omega$ only through the parity vector of its orbit multiplicities:
\[
  \Grundy(R_T(\Omega))
  =\bigoplus_{[K]}(m_K\bmod 2)\,\Grundy(R(G,K,T)),
  \qquad
  t_K:=\Grundy(R(G,K,T)).
\]
Geometry therefore enters only by choosing orbit multiplicities; the abstract pair $(G,T)$ chooses
a finite table of transitive-template nimbers. The three applications compute that table:

\begin{enumerate}[label=(\roman*)]
  \item \textbf{Two selected points} (Section~\ref{sec:pairs}) induce a dihedral group $D_{2m}$ for
    every $m$ and either parity of the rotation order. The templates are exactly a cycle, a path,
    or an empty graph, so the nimbers are the classical Dawson's-chess values and are unbounded.
  \item \textbf{Three selected points} (Section~\ref{sec:triples}) must contain the central
    involution, which forces the even-rotation groups $D_{4n}$ together with the Klein-four
    boundary. The templates are exactly an empty graph, $K_4$, a Möbius ladder, a prism, or an
    ordinary ladder, all of value in $\{0,1\}$.
  \item \textbf{Polyhedral triples} (Section~\ref{sec:polyhedral}) are the remaining tame
    possibilities, $S_4$ and $A_5$; $A_4$ cannot occur. This is a finite boundary and is settled by
    exact computation of the templates.
\end{enumerate}

Section~\ref{sec:arithmetic} converts these tables into closed field formulas, P/N congruences, and
a density theorem. Section~\ref{sec:boundaries} collects the converse realization statements and
states the boundaries of the classification.

\subsection{Conventions and standing hypotheses}

Throughout, $q$ is odd, $C\subset\PG(2,q)$ is a nonsingular conic, and
\[
  p=\operatorname{char}\F_q .
\]
We work in the \emph{tame} case, meaning that $p$ does not divide the order of the subgroup under
consideration. We write
\[
  D_{2m}=\langle r,s : r^m=s^2=1,\ srs=r^{-1}\rangle,
\]
so that $D_{2m}$ has order $2m$. The pair family below is indexed by $D_{2m}$ and the triple family
by $D_{4n}$ with $m=2n$; the two indices are kept distinct at every theorem boundary.

All game values refer to the continuation in which only conic points may be played. They need not
equal the value of a larger position in which further off-conic moves remain available.

\subsection{Main theorem}

\begin{theorem}[Classification, tame dihedral configurations]
\label{thm:main}
Let $q$ be odd and let $S$ be a legal set of selected off-conic points whose induced involutions
generate a tame dihedral subgroup $G\le\PGL_2(q)$. Then $|S|\in\{2,3\}$, and:
\begin{enumerate}[label=(\alph*)]
  \item every transitive residual template of $(G,T)$ is an empty graph, $K_4$, a cycle $\Cyc_{2m}$,
    a path $P_m$, an ordinary ladder $L_k$, a prism $\Cyc_{2n}\square K_2$, or a Möbius ladder
    $\Mob_{4n}$;
  \item the conic-only residual Grundy value is the xor of those template nimbers weighted by the
    parity of their orbit multiplicities, and is an explicit congruence function of $q$, periodic
    in $q$ within each torus family;
  \item every template in (a) is realized over infinitely many prime fields.
\end{enumerate}
The two-point family realizes every $m$ and has unbounded values; the three-point family realizes
exactly the even rotation orders and has values in $\{0,1\}$.
\end{theorem}

The statement is proved in pieces: (a) in Sections~\ref{sec:pairs} and~\ref{sec:triples}, (b) in
Section~\ref{sec:arithmetic}, and (c) in Section~\ref{sec:boundaries}.

\phasenote{C307}{Theorem~\ref{thm:main} is stated here at the strength the migrated source
manuscript proves. C307 extends it to the polyhedral families using the C284 completeness theorem,
the C289 $\rho$ refinement, and the C290 congruence laws, and must name the three exclusions
(off-conic continuation, wild characteristic, growing full/subfield $\PSL_2/\PGL_2$ groups) in the
same breath, per Fable's answer~(10).}

\section{From conic saturation to Node Kayles}
\label{sec:reduction}

Choose coordinates in which
\[
  C:\ XZ=Y^2,
\]
and parameterize
\[
  c(t)=[t^2:t:1],\qquad t\in\F_q,
\]
with $c(\infty)=[1:0:0]$.

For an off-conic point $x$, projection from $x$ exchanges the two conic points on each secant
through $x$. Denote the resulting involution of $C$ by $\sigma_x$. For $x=[a:b:c]$ this involution
is represented by the traceless matrix
\[
  M_x=\begin{pmatrix} b&-a\\ c&-b\end{pmatrix},
\]
acting on $\PP^1(q)$ as
\begin{equation}
  \sigma_x(t)=\frac{bt-a}{ct-b}.
  \label{eq:involution}
\end{equation}

Let $S$ be a legal set of selected off-conic points. A conic point is \emph{dead} when it lies on a
line through two points of $S$.

\begin{lemma}[Pair-product fixed points are exactly the dead conic points]
\label{lem:dead}
For distinct selected points $x,y$ and $u\in C$,
\begin{equation}
  u\in xy \quad\Longleftrightarrow\quad (\sigma_x\sigma_y)(u)=u .
  \label{eq:dead}
\end{equation}
\end{lemma}

\begin{proof}
Let $v=\sigma_y(u)$. The points $u,v,y$ are collinear. If $\sigma_x(v)=u$, then $u,v,x$ are also
collinear. When $u\ne v$, the line through $u,v$ is unique, so it contains $x$ and $y$, and hence
$u\in xy$. When $u=v$, both $x$ and $y$ lie on the unique tangent to $C$ at $u$, so again $u\in xy$.

Conversely, suppose $u\in xy$. Let $v$ be the other point of $C\cap xy$, counted with multiplicity.
Projection from $y$ exchanges $u$ and $v$, and projection from $x$ exchanges $v$ and $u$. Thus
$\sigma_x\sigma_y(u)=u$.
\end{proof}

Define
\begin{equation}
  D_S=\bigcup_{\{x,y\}\subset S}\Fix(\sigma_x\sigma_y).
  \label{eq:DS}
\end{equation}
On the live vertex set $C\setminus D_S$, join $u$ to $\sigma_x(u)$ for every $x\in S$, suppressing
loops and repeated edges. Call the resulting graph $R_S$.

\begin{theorem}[Exact residual reduction]
\label{thm:residual}
A legal conic move at $u\in C\setminus D_S$ deletes exactly the closed neighbourhood
$N_{R_S}[u]$. Hence play restricted to the conic is Node Kayles on $R_S$.
\end{theorem}

\begin{proof}
After selecting $u$, a live conic point $v\ne u$ becomes dead exactly when $u,v,x$ are collinear for
some $x\in S$. By definition of the projection involution this is equivalent to $v=\sigma_x(u)$.
Thus the newly unavailable conic points are precisely $u$ and its graph neighbours.
\end{proof}

The graph is static throughout the conic-only continuation. A line through two selected conic points
contains no third conic point, while a line through a newly selected conic point and an old
off-conic point $x\in S$ removes exactly its projection partner. Later moves therefore pass to
induced subgraphs of the original $R_S$ and never add edges. This is what makes the reduction
genuine Node Kayles rather than a dynamically changing graph game.

\section{Orbit templates and mod-two cancellation}
\label{sec:templates}

Let a finite group $G$ act on a finite set $\Omega$, and let $T$ be a finite set of involutions
generating $G$. Define
\begin{equation}
  D_T(\Omega)=\bigcup_{\{\sigma,\tau\}\in\binom{T}{2}}\Fix_\Omega(\sigma\tau),
  \label{eq:DT}
\end{equation}
where the unordered-pair notation is unambiguous because $\tau\sigma=(\sigma\tau)^{-1}$ has the same
fixed points as $\sigma\tau$. Equivalently, for involutions,
\[
  \sigma\tau(x)=x\quad\Longleftrightarrow\quad \sigma(x)=\tau(x),
\]
which is the common-projection-partner condition of Lemma~\ref{lem:dead}. Define also
\begin{equation}
  R_T(\Omega)=\Sch(G,\Omega,T)\bigl[\Omega\setminus D_T(\Omega)\bigr],
  \label{eq:RT}
\end{equation}
with loops and repeated edges suppressed. For a subgroup $K\le G$ write
\begin{equation}
  R(G,K,T)=R_T(G/K).
  \label{eq:template}
\end{equation}

\begin{theorem}[Orbit-template reduction]
\label{thm:orbit-template}
If $\Omega\cong\bigsqcup_{[K]}m_K(G/K)$, then Sprague-Grundy additivity gives
\begin{equation}
  \Grundy(R_T(\Omega))=\bigoplus_{[K]}(m_K\bmod 2)\,\Grundy(R(G,K,T)).
  \label{eq:orbit-template}
\end{equation}
\end{theorem}

\begin{proof}
Every generator preserves each $G$-orbit, so both the Schreier edges and the pair-product fixed
points split orbit by orbit. Although $D_T(\Omega)$ need not itself be $G$-stable when $T$ is not
closed under conjugation, every $G$-equivariant bijection commutes with every element of $T$; it
therefore transports $D_T$ and the residual graph. Thus each orbit of type $G/K$ contributes the
well-defined template $R(G,K,T)$. Grundy values xor over disjoint graph components, and two equal
nimbers cancel, giving~\eqref{eq:orbit-template}.
\end{proof}

Consequently, once the finite table $t_K=\Grundy(R(G,K,T))$ is known, the ambient action enters only
through the parity vector $(m_K\bmod 2)_{[K]}$. The rest of the paper computes these tables.

\subsection{Burnside reformulation}

Fix $G$ and $T$ as above and define $\Phi_T(\Omega)=\Grundy(R_T(\Omega))$ for every finite $G$-set
$\Omega$.

\begin{proposition}
\label{prop:burnside}
The map $\Phi_T$ is additive under disjoint union and extends uniquely to a homomorphism from the
additive group underlying the Burnside ring,
\begin{equation}
  \Phi_T:A(G)\longrightarrow(\mathbb N_0,\oplus).
  \label{eq:phi}
\end{equation}
It vanishes on $2A(G)$ and therefore factors through the mod-two Burnside group
$A(G)\otimes_{\Z}\F_2$.
\end{proposition}

\begin{proof}
The construction $R_T$ respects disjoint unions, and Grundy values xor over graph components, so
$\Phi_T(\Omega\sqcup\Omega')=\Phi_T(\Omega)\oplus\Phi_T(\Omega')$. The Grothendieck-group universal
property gives the extension. Finally $\Phi_T(\Omega\sqcup\Omega)=0$, so doubled classes lie in the
kernel.
\end{proof}

\begin{corollary}[Bulk cancellation]
\label{cor:bulk}
If a finite $G$-set has $f$ free orbits and exceptional part $\Omega_{\mathrm{exc}}$, then
\begin{equation}
  \Phi_T(\Omega)=(f\bmod 2)\,\Phi_T(G/1)\oplus\Phi_T(\Omega_{\mathrm{exc}}).
  \label{eq:bulk}
\end{equation}
Arbitrarily many generic free orbits therefore contribute either nothing or a single parity bit.
\end{corollary}

\begin{remark}
\label{rem:burnside-scope}
This is a reformulation of Theorem~\ref{thm:orbit-template}, not an additional source of
computational information: the additive group of $A(G)$ is already free on the transitive $G$-sets,
and no compatibility with Burnside-ring multiplication is asserted.
\end{remark}

\phasenote{C308}{Proposition~\ref{prop:burnside} and Corollary~\ref{cor:bulk} are the
Lean-anchored results of this section (C262, \texttt{Burnside.lean}). C308 attaches the formal
anchor citation to the corollary, per Fable's answer~(3).}

\section{Two selected points: every dihedral order}
\label{sec:pairs}

\subsection{Legal dihedral configurations}
\label{subsec:legal}

\begin{lemma}[Collinearity of reflection centres]
\label{lem:collinear}
For any dihedral subgroup of $\PGL_2(q)$, regardless of the parity of its rotation order, all
reflection centres lie on one projective line.
\end{lemma}

\begin{proof}
Over the splitting field a dihedral normal form represents every reflection by a matrix proportional
to
\[
  \begin{pmatrix} 0&\lambda\\ 1&0\end{pmatrix},
\]
whose centre has coordinate $b=0$. The line spanned by the reflection centres is Galois invariant
and therefore descends to $\F_q$.
\end{proof}

A legal selected set may have any size, so we record the complete taxonomy of legal sets whose
induced involutions generate a dihedral group. By Lemma~\ref{lem:collinear} a legal set contains at
most two reflections, and it contains at most one central involution $z$. Hence a legal dihedral
configuration has either two or three selected points:

\begin{table}[H]
\centering
\begin{tabular}{lll}
\toprule
selected involutions & generated group & admissible $m$\\
\midrule
two reflections & $D_{2m}$ & every $m\ge3$, both parities\\
one reflection and the central $z$ & $V_4$ & $m=2$ boundary\\
two reflections and the central $z$ & $D_{4n}$, $m=2n$ & even rotation order only\\
\bottomrule
\end{tabular}
\caption{Legal dihedral configurations. A set of four or more selected points cannot induce a
dihedral group, since it would contain three collinear reflection centres.}
\label{tab:legal}
\end{table}

The third row is the group $D_{4n}$ analysed in Section~\ref{sec:triples}. The first row is the
general dihedral family: any two selected off-conic points $x,y$ induce involutions
$\sigma_x,\sigma_y$ generating
\[
  G=\langle\sigma_x,\sigma_y\rangle\cong D_{2m},\qquad m=\ord(\sigma_x\sigma_y),
\]
with no central involution required, so this family exists for \emph{every} $m$, including the odd
rotation orders that admit no legal triple.

Throughout this section $\{x,y\}$ is a legal pair, $G\cong D_{2m}$ is tame ($p\nmid 2m$), and the
generating involution set is
\[
  T=\{\sigma_x,\sigma_y\}=\{s,sr^d\},\qquad \gcd(d,m)=1 .
\]
Since $\gcd(d,m)=1$ we may replace $r$ by $r^d$ and take $d=1$, so $T=\{s,sr\}$ and
$\sigma_x\sigma_y=r$.

\subsection{Point stabilizers and templates}

As in Proposition~\ref{prop:stabilizers} below, every tame point stabilizer of $G$ on $\PP^1(q)$ is
cyclic, hence conjugate to $1$, to the rotation subgroup $\langle r\rangle$, or to a reflection
subgroup $\langle sr^e\rangle$. The reflections form one conjugacy class when $m$ is odd and two
classes, represented by $\langle s\rangle$ and $\langle sr\rangle$, when $m$ is even.

\begin{theorem}[Two-generator templates]
\label{thm:pair-templates}
For $T=\{s,sr\}$,
\begin{equation}
  R(D_{2m},1,T)\cong \Cyc_{2m},\qquad
  R(D_{2m},\langle sr^e\rangle,T)\cong P_m,\qquad
  R(D_{2m},\langle r\rangle,T)=\varnothing,
  \label{eq:pair-templates}
\end{equation}
where $\Cyc_{2m}$ is the cycle on $2m$ vertices and $P_m$ the path on $m$ vertices.
\end{theorem}

\begin{proof}
\emph{Free orbit.} The regular action has no deleted vertices, so the template is the Cayley graph
$\Cay(G,T)$. It is $2$-regular, and its edges alternate between the two involution generators. Since
$T$ generates $G$, that alternation visits all $2m$ group elements in a single closed walk. Hence it
is one $2m$-cycle.

\emph{Rotation orbit.} The two cosets of $\langle r\rangle$ are each fixed by
$\sigma_x\sigma_y=r\in\langle r\rangle$, so both are deleted by~\eqref{eq:DT} and the template is
empty.

\emph{Reflection orbit.} Label the cosets $D_{2m}/\langle s\rangle$ by $j\in\Z/m$ via
$j\leftrightarrow r^j\langle s\rangle$. Left multiplication gives
\begin{equation}
  s:j\mapsto-j,\qquad sr:j\mapsto-1-j,\qquad\text{product } r:j\mapsto j+1 .
  \label{eq:pair-coset-action}
\end{equation}
The deleted vertices are $\Fix(\sigma_x\sigma_y)=\Fix(r)$ on this coset set; but $r:j\mapsto j+1$ is
fixed-point-free on $\Z/m$, so no vertex of a reflection orbit is deleted. The two reflection maps
in~\eqref{eq:pair-coset-action} have product the $m$-cycle $r$, so their edge union is connected of
maximum degree two. Their fixed points, which produce suppressed loops rather than edges, are the
solutions of $2j\equiv0$ and $2j\equiv-1\pmod m$; counting solutions in each parity of $m$ leaves
exactly two vertices of degree one and $m-2$ of degree two. A connected graph of maximum degree two
with exactly two degree-one vertices is the path $P_m$. For a general reflection stabilizer
$\langle sr^e\rangle$, the automorphism $r\mapsto r,\ s\mapsto sr^e$ of $D_{2m}$ carries
$\langle s\rangle$ to $\langle sr^e\rangle$ and $T=\{s,sr\}$ to another two-reflection set with
product $r$, so the same computation applies and the template is $P_m$ independently of $e$.
\end{proof}

\subsection{Template nimbers}

Node Kayles on a path is Dawson's chess: deleting a closed neighbourhood removes an interior vertex
together with its at most two neighbours, splitting $P_m$ into two shorter paths, which is the move
rule of the octal game $0.137$ \cite{WinningWays}. Its Grundy sequence
\begin{equation}
  \Grundy(P_m)=\mathrm{A}002187(m)=0,1,1,2,0,3,1,1,0,3,3,2,2,4,0,5,\dots\qquad(m=0,1,2,\dots)
  \label{eq:dawson}
\end{equation}
is eventually periodic with period $34$ \cite{OEIS-A002187}.

For the cycle, every vertex of $\Cyc_k$ has the same closed neighbourhood shape, and deleting it
removes three consecutive vertices, leaving $P_{k-3}$. Vertex-transitivity makes this the only
option up to isomorphism, so
\begin{equation}
  \Grundy(\Cyc_k)=\mex\{\Grundy(P_{k-3})\}\qquad(k\ge3).
  \label{eq:cycle-mex}
\end{equation}

\begin{lemma}[Free templates are neutral]
\label{lem:cycle-zero}
For every $m\ge2$, $\Grundy(\Cyc_{2m})=0$.
\end{lemma}

\begin{proof}
By~\eqref{eq:cycle-mex}, $\Grundy(\Cyc_{2m})=\mex\{\Grundy(P_{2m-3})\}$, which is zero unless
$\Grundy(P_{2m-3})=0$. The index $2m-3$ is odd, so it suffices to know that the Dawson
sequence~\eqref{eq:dawson} has no zero at an odd index. Its period is $34$, which is even, so the
index parity is constant along each residue class modulo the period; hence it suffices to check the
finitely many indices up to the end of the first full period together with the pre-periodic part,
and there every zero of $\mathrm{A}002187$ occurs at an even index. Therefore
$\Grundy(P_{2m-3})\ne0$ and the mex is zero.
\end{proof}

Free orbits therefore never contribute to the value of a two-point residual.

\begin{theorem}[Pair orbit formula]
\label{thm:pair-orbit}
Let $\Omega$ be a tame $D_{2m}$-set with $f$ free orbits and $\rho$ reflection orbits. Then
\begin{equation}
  \Grundy(R_T(\Omega))
  =(f\bmod2)\,\Grundy(\Cyc_{2m})\oplus(\rho\bmod2)\,\Grundy(P_m)
  =(\rho\bmod2)\,\Grundy(P_m).
  \label{eq:pair-orbit}
\end{equation}
\end{theorem}

\begin{proof}
Theorem~\ref{thm:orbit-template} gives the first equality, with the rotation orbits contributing the
empty template. Lemma~\ref{lem:cycle-zero} kills the free term.
\end{proof}

\begin{corollary}[Period $34$ in the rotation order]
\label{cor:period34}
For a fixed reflection-orbit parity, the two-point residual value $(\rho\bmod2)\,\Grundy(P_m)$ is
eventually periodic in $m$ with period $34$.
\end{corollary}

\begin{proof}
Immediate from~\eqref{eq:pair-orbit} and the eventual period-$34$ periodicity
of~\eqref{eq:dawson}.
\end{proof}

The Dawson zeros $\Grundy(P_m)=0$ at even $m\in\{4,8,14,20,24,28,34,\dots\}$ give infinite families
of even-order dihedral groups whose two-point game is always a P-position. This is the
two-generator analogue of the always-deleted rotation torus, but arises from a nontrivial path
cancellation rather than from a deletion.

\section{Three selected points: the dihedral ladder family}
\label{sec:triples}

\subsection{The Klein-four boundary}

The smallest dihedral boundary is $D_4\cong V_4$. A generating involution triple is necessarily the
set of all three nonidentity elements.

\begin{theorem}[Klein-four boundary]
\label{thm:v4}
Let a legal tame conic-involution triple generate $G\cong V_4$, and let $s$ be the number of its
three involutions having rational fixed points on $C$. The realizable split counts satisfy
\begin{equation}
  s\in\{1,3\}\ (q\equiv1\bmod4),
  \qquad
  s\in\{0,2\}\ (q\equiv3\bmod4).
  \label{eq:v4-split}
\end{equation}
The dead set has size $2s$, and the residual graph is
\begin{equation}
  \frac{q+1-2s}{4}\,K_4 ,
  \label{eq:v4-graph}
\end{equation}
a coefficient denoting a disjoint union of that many copies. Consequently
\begin{equation}
  \Grundy(R_T)=\frac{q+1-2s}{4}\bmod2 .
  \label{eq:v4-value}
\end{equation}
\end{theorem}

\begin{proof}
The product of any two nonidentity elements of $V_4$ is the third, so~\eqref{eq:DT} deletes exactly
the union of the fixed-point sets of the three involutions. A tame point stabilizer in $\PGL_2(q)$
is cyclic, so two distinct nonidentity elements of $V_4$ cannot fix the same conic point. Each split
involution has two rational fixed points and each nonsplit involution has none, giving $|D_T|=2s$.

The split-count restriction is the parity law for a tame Klein four subgroup of $\PGL_2(q)$: the
number of split nonidentity elements has parity opposite to $(q-1)/2$, giving exactly the
alternatives in~\eqref{eq:v4-split}. Equivalently the orbit equation $q+1=2s+4f$ forces the same
parity, and the standard split/nonsplit normal forms realize both listed values. In particular
divisibility of $q+1-2s$ by four is automatic.

The $V_4$-action on the live points is free: a nontrivial stabilizer would again contain one of the
involutions whose fixed points were deleted. Every live orbit therefore has four points, and on such
an orbit the three involutions give the three perfect matchings of $K_4$. Since one Node-Kayles move
deletes a whole $K_4$, each orbit has Grundy value one, and disjoint-union additivity
gives~\eqref{eq:v4-value}.
\end{proof}

In the graph notation below this is the $n=1$ boundary $\Mob_4=K_4$, while the reflection-stabilized
template is $L_0=\varnothing$. It is treated separately because the $T_d$ parametrization and the
pendant-ladder proof for $n\ge2$ degenerate at $n=1$.

\subsection{The legal triples \texorpdfstring{$T_d$}{Td}}

No dihedral group whose rotation subgroup has odd order can arise from a legal generating involution
triple: all of its involutions are reflections, and their centres are collinear by
Lemma~\ref{lem:collinear}. For the remainder of the section fix $n\ge2$ and suppose
\[
  G\cong D_{4n}=\langle r,s: r^{2n}=s^2=1,\ srs=r^{-1}\rangle,
  \qquad z=r^n .
\]
The noncentral involutions are the reflections $sr^e$, so three noncentral involution centres are
collinear and cannot be a legal selected triple. A legal generating involution triple must therefore
contain $z$ and two reflections. After rechoosing the reflection generator and interchanging the
reflections, it has the form
\begin{equation}
  T_d=\{z,s,sr^d\}.
  \label{eq:Td}
\end{equation}

\begin{proposition}
\label{prop:Td-generates}
The triple $T_d$ generates $D_{4n}$ exactly when $\gcd(d,n)=1$.
\end{proposition}

\begin{proof}
The subgroup contains $s(sr^d)=r^d$ and $r^n$, so its rotation subgroup is
$\langle r^d,r^n\rangle=\langle r^{\gcd(d,n)}\rangle$. This equals $\langle r\rangle$ precisely when
$\gcd(d,n)=1$.
\end{proof}

Under an automorphism $r\mapsto r^u$, $s\mapsto r^v s$ with $u\in(\Z/2n)^\times$, the parameter $d$
changes to $\pm ud$. Therefore if $n$ is even every admissible $d$ is odd and there is one
automorphism class, while if $n$ is odd there are two classes, represented by odd and even $d$. The
central centre is not on the reflection-centre line, so every generating triple described above is
geometrically legal.

\subsection{Free-orbit templates}

Let $\Mob_{4n}$ denote the graph obtained from $\Cyc_{4n}$ by joining antipodal vertices, and let
$L_k=P_k\square K_2$ denote the ladder with $k$ rungs.

\begin{theorem}
\label{thm:free-template}
For the regular action,
\begin{equation}
  R(D_{4n},1,T_d)\cong
  \begin{cases}
    \Mob_{4n}, & d\text{ odd},\\
    \Cyc_{2n}\square K_2, & d\text{ even}.
  \end{cases}
  \label{eq:free-template}
\end{equation}
The even case occurs only when $n$ is odd.
\end{theorem}

\begin{proof}
No nonidentity element fixes a point in the regular action, so no vertices are deleted.

The two reflection generators alternate along the cosets of $\langle r^d\rangle$. If $d$ is odd then
$\gcd(d,2n)=1$, so their edges form one $4n$-cycle. Multiplication by $z=r^n$ joins opposite
vertices: centrality of $z$ makes this chord matching invariant under the reflection cycle, and
after two reflection steps one advances by $d$, so $z$ advances by exactly half of that $4n$-cycle.
This gives $\Mob_{4n}$.

If $d$ is even then $n$ is odd and $\gcd(d,2n)=2$. The reflection edges form two $2n$-cycles, and
$z$ joins corresponding vertices of the two cycles. This is $\Cyc_{2n}\square K_2$.
\end{proof}

\begin{theorem}
\label{thm:mobius-value}
For every $n\ge2$, $\Grundy(\Mob_{4n})=1$.
\end{theorem}

\begin{proof}
The graph is vertex-transitive. After a move at $v$, the graph $\Mob_{4n}-N[v]$ is the
opposite-end-pendant ladder on $2n-3$ rungs. To see this explicitly, label $\Mob_{4n}$ by $\Z/(4n)$,
with cycle edges $i\sim i+1$ and chords $i\sim i+2n$, and take $v=0$. Removing
$N[0]=\{0,1,2n,4n-1\}$ leaves the paired vertices
\[
  2+i\longleftrightarrow 2n+2+i,\qquad 0\le i\le 2n-4,
\]
as the $2n-3$ rungs. The unpaired vertex $2n-1$ is pendant at the right end of the first rail, while
$2n+1$ is pendant at the left end of the second rail. This is exactly Brown et al.'s third ladder
variation ${}_{-}L^-_{(2n-3)\times2}$ \cite{BrownEtAl2020}.

Their ladder theorem states $\Grundy({}_{-}L^-_{k\times2})=0$ for every integer $k\ge1$. Here
$k=2n-3$ is odd and at least one. Their $k=1$ base graph is $P_4$, with value zero, and $k=3$ is
already covered by the same theorem; the first two values arising here are therefore included
explicitly, not inferred from an eventual pattern. Hence every option of $\Mob_{4n}$ has value zero,
and $\Grundy(\Mob_{4n})=\mex\{0\}=1$.
\end{proof}

\begin{theorem}
\label{thm:prism-value}
For every $k\ge3$, $\Grundy(\Cyc_k\square K_2)=0$.
\end{theorem}

Brown et al.'s prism corollary \cite{BrownEtAl2020} states this for every $k\ge3$, with no parity or
residue restriction. In particular it applies to the only prisms used here, where $k=2n\ge6$ and $n$
is odd.

\subsection{Point stabilizers in the tame projective action}

Let $A=\langle r\rangle\cong \Z/2n$ and $K_e=\langle sr^e\rangle$.

\begin{proposition}
\label{prop:stabilizers}
Every point stabilizer of $D_{4n}$ on $\PP^1(q)$ is conjugate to one of
\begin{equation}
  1,\qquad A,\qquad K_0,\qquad K_1 .
  \label{eq:stabilizers}
\end{equation}
\end{proposition}

\begin{proof}
A point stabilizer in $\PGL_2(q)$ is contained in a Borel subgroup. Its normal unipotent subgroup has
order a power of $p$; since $p\nmid|G|$, the intersection of $G$ with that subgroup is trivial.
Projection to the cyclic quotient of the Borel by its unipotent subgroup is therefore injective on
the point stabilizer, so the stabilizer is cyclic.

If the stabilizer contains a nonidentity rotation, the point is one of the two fixed points of the
rotation torus over the splitting field and is therefore fixed by all of $A$. An apparent
intermediate rotation stabilizer $\langle r^k\rangle$ immediately enlarges to $A$, and it cannot
enlarge further inside $D_{4n}$, since adjoining any reflection to $A$ would make the stabilizer
dihedral and hence noncyclic. Otherwise the stabilizer is generated by a reflection. Because the
rotation order is even, the reflections form two conjugacy classes, represented by $K_0$ and $K_1$.
\end{proof}

\begin{proposition}
\label{prop:rotation-empty}
The rotation-stabilized template is empty: $R(D_{4n},A,T_d)=\varnothing$.
\end{proposition}

\begin{proof}
The two cosets of $A$ are fixed by $s(sr^d)=r^d\in A$, so both are deleted by~\eqref{eq:DT}.
\end{proof}

\subsection{Reflection-stabilized templates}

Label the cosets $D_{4n}/K_e$ by $j\in\Z/(2n)$ via $j\leftrightarrow r^jK_e$. The generators act as
\begin{equation}
  z:j\mapsto j+n,\qquad s:j\mapsto e-j,\qquad sr^d:j\mapsto e-d-j .
  \label{eq:triple-coset-action}
\end{equation}
The two reflection maps have product $j\mapsto j+d$. The deleted vertices are the solutions of
\begin{equation}
  2j=e+n\quad\text{or}\quad 2j=e-d+n \pmod{2n}.
  \label{eq:deleted}
\end{equation}

\begin{lemma}[A double cover of a tree is untwisted]
\label{lem:double-cover}
Let a fixed-point-free involution $z$ act on a graph $X$. Suppose a $z$-invariant edge set $E_0$
descends to a graph $Y=X/\langle z\rangle$, with every edge of $Y$ having exactly two $z$-paired
lifts. If $Y$ is a tree, then $(V(X),E_0)$ is the disjoint union of two copies of $Y$, exchanged by
$z$. After adding every edge $v\sim z(v)$, the result is $Y\square K_2$.
\end{lemma}

\begin{proof}
Choose a root of $Y$ and one of its two lifts. Every vertex of a tree has a unique path from the
root, so lifting that path selects a unique lift of every vertex and produces one copy of $Y$.
Applying $z$ produces the other copy. There is no monodromy because the tree has no cycle. The added
$z$-edges join the two lifts of each quotient vertex, which are exactly the $K_2$-fibres of the
Cartesian product.
\end{proof}

\begin{theorem}
\label{thm:reflection-template}
For odd $d$,
\begin{equation}
  R(D_{4n},K_e,T_d)\cong L_{n-1},
  \label{eq:refl-odd}
\end{equation}
independently of $e$. For even $d$, necessarily $n$ odd,
\begin{equation}
  R(D_{4n},K_e,T_d)\cong
  \begin{cases}
    L_n, & e\text{ even},\\
    L_{n-2}, & e\text{ odd}.
  \end{cases}
  \label{eq:refl-even}
\end{equation}
\end{theorem}

\begin{proof}
Because $\gcd(d,n)=1$, reduction modulo $n$ makes translation by $d$ transitive on $\Z/n$,
regardless of the parity of $d$. Quotient the coset set by the fixed-point-free central involution
$z:j\mapsto j+n$. On the quotient, the two reflection maps are involutions whose product is this
transitive translation. Their simple edge union is connected and has maximum degree two. The two
quotient reflections have two fixed points in total: each has one when $n$ is odd, while for even
$n$ one has two and the other has none. Suppressing the corresponding loops leaves $n-1$ edges, so
the quotient reflection graph is the path $P_n$, with those fixed points as its endpoints.

Each solvable congruence in~\eqref{eq:deleted} has two solutions differing by $n$, hence deletes one
$z$-orbit. After reduction modulo $n$, the two congruences are exactly the fixed-point equations for
the two quotient reflections, so every deleted $z$-orbit is an endpoint of $P_n$. If $d$ is odd, the
right-hand sides in~\eqref{eq:deleted} have opposite parity, so exactly one endpoint orbit is
deleted and the live quotient is $P_{n-1}$. If $d$ is even, then $n$ is odd; for even $e$ neither
congruence is solvable and the live quotient is $P_n$, while for odd $e$ both are solvable and the
live quotient is $P_{n-2}$.

The reflection edges above quotient path edges satisfy Lemma~\ref{lem:double-cover}, and the
$z$-edges are its fibre edges. A reflection above a suppressed endpoint loop either gives loops
upstairs or duplicates its $z$-edge, neither of which changes the underlying simple graph. The three
live quotients therefore lift to $L_{n-1}$, $L_n$, and $L_{n-2}$ respectively.
\end{proof}

Brown et al.\ \cite{BrownEtAl2020} prove
\begin{equation}
  \Grundy(L_k)=
  \begin{cases}
    1, & k\text{ odd},\\
    0, & k\text{ even},
  \end{cases}
  \label{eq:ladder-value}
\end{equation}
and therefore
\begin{equation}
  \Grundy(R(D_{4n},K_e,T_d))=
  \begin{cases}
    \mathbf 1_{2\mid n}, & d\text{ odd},\\
    1, & d\text{ even}.
  \end{cases}
  \label{eq:refl-value}
\end{equation}

\subsection{Complete orbit formulas}

Let $\Omega$ be a finite $D_{4n}$-set all of whose point stabilizers are conjugate to one of
$1,A,K_0,K_1$; by Proposition~\ref{prop:stabilizers} the natural tame action on $\PP^1(q)$ has this
property. Write
\begin{equation}
  \Omega\cong f(D_{4n}/1)\sqcup b_A(D_{4n}/A)\sqcup b_0(D_{4n}/K_0)\sqcup b_1(D_{4n}/K_1)
  \label{eq:multiplicities}
\end{equation}
with nonnegative integer multiplicities.

\begin{theorem}[Triple orbit formula]
\label{thm:triple-orbit}
For odd $d$,
\begin{equation}
  \Grundy(R_{T_d}(\Omega))=(f\bmod2)\oplus\bigl(\mathbf 1_{2\mid n}((b_0+b_1)\bmod2)\bigr).
  \label{eq:triple-odd}
\end{equation}
For even $d$, necessarily $n$ odd,
\begin{equation}
  \Grundy(R_{T_d}(\Omega))=(b_0+b_1)\bmod2 .
  \label{eq:triple-even}
\end{equation}
The multiplicity $b_A$ does not occur because its template is empty.
\end{theorem}

For the natural action on $\PP^1(q)$, each exceptional class occurs at most once and the exceptional
point sets are pairwise disjoint: otherwise one point stabilizer would contain either two distinct
reflections or a reflection and a nonidentity rotation, contradicting the cyclicity conclusion of
Proposition~\ref{prop:stabilizers}. (Here $2n\ge4$, so the two reflection classes and the central
rotation are genuinely distinct.) Writing the orbit indicators as
$\varepsilon,a_0,a_1\in\{0,1\}$,
\begin{equation}
  q+1=4nf+2\varepsilon+2n(a_0+a_1).
  \label{eq:orbit-equation}
\end{equation}

The complete transitive-template table for the dihedral configurations, including the Klein-four
boundary, is therefore:

\begin{table}[H]
\centering
\begin{tabular}{llcc}
\toprule
group and triple & stabilizer & residual template & $\Grundy$\\
\midrule
$V_4$, all three involutions & $1$ & $K_4$ & $1$\\
$V_4$, all three involutions & nontrivial cyclic & $\varnothing$ & $0$\\
$D_{4n}$, odd $d$ & $1$ & $\Mob_{4n}$ & $1$\\
$D_{4n}$, even $d$ & $1$ & $\Cyc_{2n}\square K_2$ & $0$\\
$D_{4n}$, any $d$ & $A$ & $\varnothing$ & $0$\\
$D_{4n}$, odd $d$ & $K_e$ & $L_{n-1}$ & $\mathbf 1_{2\mid n}$\\
$D_{4n}$, even $d$, even $e$ & $K_e$ & $L_n$ & $1$\\
$D_{4n}$, even $d$, odd $e$ & $K_e$ & $L_{n-2}$ & $1$\\
\midrule
$D_{2m}$, two reflections & $1$ & $\Cyc_{2m}$ & $0$\\
$D_{2m}$, two reflections & $\langle sr^e\rangle$ & $P_m$ & $\mathrm{A}002187(m)$\\
$D_{2m}$, two reflections & $\langle r\rangle$ & $\varnothing$ & $0$\\
\bottomrule
\end{tabular}
\caption{The complete transitive-template table for legal tame dihedral configurations. The graph
families of Theorem~\ref{thm:main} are a corollary of this one table; the ambient projective line
merely supplies its orbit multiplicities.}
\label{tab:templates}
\end{table}

\section{Polyhedral triples: the finite boundary}
\label{sec:polyhedral}

Beyond the dihedral configurations the tame possibilities are finite. The next subgroup in the usual
polyhedral list, $A_4$, cannot occur as a generated group here: its only involutions are the three
nonidentity elements of its normal Klein four subgroup, so any subgroup of $A_4$ generated by
involutions is contained in that $V_4$, not all of $A_4$. The remaining polyhedral possibilities are
$S_4$ and $A_5$.

\phasenote{C307}{This section is the structural placeholder for the promoted polyhedral content.
Per Fable's change~4 it must \emph{lead with completeness} --- the C284 classification theorem that
all admissible $S_4/A_5$ coset templates are tabulated, with $A_4$ impossible --- and then present
the $A_5$ $(3,5,5;\rho=3/5)$ split as the load-bearing refinement, proved as an $\Aut(G)$-orbit
invariant (C289) with no stronger geometric interpretation. The ten-row
$(\sigma,\rho,\text{class size},t_1\ldots t_5)$ table belongs here in the body; edge lists and
replay counts go to Appendix~\ref{app:evidence} (Fable's answer~4). C289's mirror lemma upgrades six
$t_1=0$ entries from computed to proved; the $\rho=3$ regular zero stays computational and must be
visibly marked as such. The free-orbit table below is the only polyhedral material the source
manuscript carries and is retained here unchanged pending that integration.}

\subsection{Regular polyhedral templates}

For a regular action there are no deleted vertices, so the template is the cubic Cayley graph
$\Cay(G,T)$. Generating involution triples are classified by the sorted orders of their three
pairwise products. Exact Node-Kayles computation gives Table~\ref{tab:polyhedral-regular}.

\begin{table}[H]
\centering
\begin{tabular}{llcc}
\toprule
$G$ & pair-product-order class & graph & $\Grundy$\\
\midrule
$S_4$ & $(2,3,3)$ & 24-vertex cubic & $0$\\
$S_4$ & $(2,3,4)$ & 24-vertex cubic & $0$\\
$S_4$ & $(3,3,3)$ & 24-vertex cubic & $0$\\
$S_4$ & $(3,4,4)$ & 24-vertex cubic & $0$\\
$A_5$ & $(2,3,5)$ & 60-vertex cubic, girth 4 & $1$\\
$A_5$ & $(2,5,5)$ & 60-vertex cubic, girth 4 & $1$\\
$A_5$ & $(3,3,5)$ & 60-vertex cubic, girth 6 & $0$\\
$A_5$ & $(3,5,5)$ & 60-vertex cubic, girth 6 & $0$\\
$A_5$ & $(5,5,5)$ & 60-vertex cubic, girth 9 & $0$\\
\bottomrule
\end{tabular}
\caption{Free-orbit templates and their Node-Kayles values for the polyhedral groups. These are
free-orbit values only; the computation and its cross-checks are described in
Appendix~\ref{app:regular-computation}.}
\label{tab:polyhedral-regular}
\end{table}

\begin{remark}
\label{rem:a5-commuting}
In $A_5$ the two value-one classes are exactly those containing a commuting generator pair,
equivalently a pair product of order two. This is an observation about the displayed table, not a
criterion used in the computation and not a general theorem: the analogous $(2,3,3)$ and $(2,3,4)$
classes of $S_4$ both have value zero.
\end{remark}

Together with the Klein-four and dihedral results above, Table~\ref{tab:polyhedral-regular}
completes the free-orbit row $t_1$ for every tame small-subgroup type generated by a legal
involution triple.

\section{Arithmetic synthesis}
\label{sec:arithmetic}

This section converts the template tables of Sections~\ref{sec:pairs}--\ref{sec:polyhedral} into
closed formulas over $\F_q$. Throughout, $p\nmid|G|$.

\subsection{Closed values for the three-point family}

Assume $p\nmid 4n$.

\subsubsection*{Split torus}

Suppose $2n\mid q-1$ and choose $r:t\mapsto\omega t$, $s:t\mapsto t^{-1}$ with $\omega$ of order
$2n$. The rotation torus fixes $0,\infty$, so $\varepsilon=1$. The two reflection classes are
represented by $t\mapsto t^{-1}$ and $t\mapsto\omega t^{-1}$. The first is split; the second is
split exactly when $\omega$ is a square, equivalently when $4n\mid q-1$. Put
\[
  h_-=\frac{q-1}{2n},\qquad \delta_-=\mathbf 1_{4n\mid q-1}.
\]
Then
\begin{equation}
  a_0=1,\qquad a_1=\delta_-,\qquad f=\frac{h_--1-\delta_-}{2}.
  \label{eq:split-orbits}
\end{equation}
For odd $d$,
\begin{equation}
  \Grundy(R_{T_d})
  =\left(\frac{h_--1-\delta_-}{2}\bmod 2\right)\oplus\bigl(\mathbf 1_{2\mid n}(1-\delta_-)\bigr),
  \label{eq:split-odd}
\end{equation}
and for even $d$,
\begin{equation}
  \Grundy(R_{T_d})=1-\delta_- .
  \label{eq:split-even}
\end{equation}

\subsubsection*{Nonsplit torus}

Suppose $2n\mid q+1$. Model the projective line by the norm-one group
$U=\{x\in\F_{q^2}^\times : x^{q+1}=1\}$ and choose $r:x\mapsto\omega x$, $s:x\mapsto x^{-1}$ with
$\omega\in U$ of order $2n$. The rotation torus has no rational fixed points, so $\varepsilon=0$.
The two reflection classes have fixed-point equations $x^2=1$ and $x^2=\omega$, so the first is
split and the second is split exactly when $4n\mid q+1$. Put
\[
  h_+=\frac{q+1}{2n},\qquad \delta_+=\mathbf 1_{4n\mid q+1}.
\]
Then
\begin{equation}
  a_0=1,\qquad a_1=\delta_+,\qquad f=\frac{h_+-1-\delta_+}{2},
  \label{eq:nonsplit-orbits}
\end{equation}
and
\begin{equation}
  \Grundy(R_{T_d})
  =\left(\frac{h_+-1-\delta_+}{2}\bmod 2\right)\oplus\bigl(\mathbf 1_{2\mid n}(1-\delta_+)\bigr)
  \label{eq:nonsplit-odd}
\end{equation}
for odd $d$, while for even $d$,
\begin{equation}
  \Grundy(R_{T_d})=1-\delta_+ .
  \label{eq:nonsplit-even}
\end{equation}

\begin{corollary}[Explicit P-position congruences]
\label{cor:pn-congruences}
In either torus family write $h=(q-1)/2n$ in the split case and $h=(q+1)/2n$ in the nonsplit case.
Then the conic-only residual is a P-position exactly under the following congruences:
\begin{equation}
  \begin{array}{c|c}
    \text{triple type} & \text{P exactly when}\\
    \hline
    d\text{ odd},\ n\text{ odd} & h\equiv1,2\pmod 4\\
    d\text{ odd},\ n\text{ even} & h\equiv2,3\pmod 4\\
    d\text{ even}\ (n\text{ odd}) & h\equiv0,2\pmod 4
  \end{array}
  \label{eq:pn-table}
\end{equation}
\end{corollary}

\begin{proof}
In both torus families $\delta=1$ exactly when $h$ is even. Substitution
into~\eqref{eq:split-odd}--\eqref{eq:nonsplit-even}, followed by the four cases $h\bmod4$, gives the
table.
\end{proof}

\phasenote{C307}{\textbf{Correctness hazard 1 --- blocking.} For even $h$ a second $D_{4n}$
conjugacy class with all-nonsplit reflections ($t=0$) exists, first at $q=7$, $D_8$ acting freely
with board $\Mob_8$ and value $1$ against the boxed $0$. C281 established that
equations~\eqref{eq:split-odd}--\eqref{eq:nonsplit-even} and
Corollary~\ref{cor:pn-congruences} require a $t$-case split, and that for odd $d$ exactly one of the
two classes is an N-position. C307 must apply that correction here before any polishing; the
formulas above are migrated verbatim from the source manuscript and are known to be incomplete as
stated. The $\S14$ pair value formula survives unchanged ($t\equiv1+\delta\bmod2$).}

For the even-$d$ triple the criterion has the geometric form
\begin{equation}
  R_{T_d}\text{ is P}
  \quad\Longleftrightarrow\quad
  \text{both reflection classes are split}
  \quad\Longleftrightarrow\quad
  4n\mid q\mp1,
  \label{eq:geometric-P}
\end{equation}
where the upper sign is the split-torus family and the lower sign the nonsplit-torus family. For
fixed $n$ these formulas are periodic in $q$ with period dividing $8n$ within each torus family.

\subsection{Closed values for the two-point family}

\begin{theorem}[Closed Grundy value of the pair family]
\label{thm:pair-closed}
Let $G\cong D_{2m}$ act tamely on $\PP^1(q)$ through a split torus ($m\mid q-1$) or a nonsplit torus
($m\mid q+1$), and put $\delta=\mathbf 1_{2m\mid q\mp1}$ with the upper sign for the split family
and the lower sign for the nonsplit family. Then
\begin{equation}
  \Grundy(R_T)=
  \begin{cases}
    0, & m\text{ odd},\\
    (1-\delta)\,\Grundy(P_m), & m\text{ even}.
  \end{cases}
  \label{eq:pair-closed}
\end{equation}
\end{theorem}

\begin{proof}
By~\eqref{eq:pair-orbit} the value is $(\rho\bmod2)\,\Grundy(P_m)$, so it suffices to compute the
parity of the number $\rho$ of reflection orbits.

\emph{Odd $m$.} The reflections form a single conjugacy class, so they are all split or all
nonsplit. If nonsplit, no reflection has a rational fixed point and $\rho=0$. If split, each of the
$m$ reflections has two rational fixed points, and a conic point is fixed by at most one reflection
because a tame point stabilizer is cyclic; so the reflection fixed points form a set $F$ of exactly
$2m$ points. Every point of $F$ has stabilizer a reflection subgroup of order two, hence lies in an
orbit of size $|G|/2=m$. Therefore $F$ splits into $2m/m=2$ orbits and $\rho=2$. In both cases
$\rho$ is even and the value is $0$.

\emph{Even $m=2n$.} Now the central involution $z=r^n$ exists. It commutes with every reflection
$\sigma$, so it permutes the two-point set $\Fix(\sigma)$. The stabilizer of a reflection fixed point
is exactly $\langle\sigma\rangle$, which does not contain the rotation $z$; hence $z$ fixes neither
point of $\Fix(\sigma)$ and therefore swaps them. So for each reflection class the two rational fixed
points, when they exist, lie in a single orbit, and each split reflection class contributes exactly
one reflection orbit. In the split-torus model $r:t\mapsto\omega t$, $s:t\mapsto t^{-1}$ with
$\omega$ of order $m$, the class $\langle s\rangle$ fixes $t=\pm1$ and is always split, while
$\langle sr\rangle$ fixes the solutions of $t^2=\omega^{-1}$, which are rational exactly when
$\omega^{-1}$ is a square, that is when $2m\mid q-1$; the nonsplit-torus model gives the same
statement with $q+1$. Therefore $\rho=1+\delta$, $(\rho\bmod2)=1-\delta$, and the value is
$(1-\delta)\,\Grundy(P_m)$.
\end{proof}

The even case is the two-generator counterpart of the three-generator group $D_{4n}$ with $m=2n$;
the reflection-splitting condition $2m\mid q\mp1$ is the condition $4n\mid q\mp1$
of~\eqref{eq:geometric-P}.

\begin{corollary}[P-positions of the pair family]
\label{cor:pair-P}
The two-point residual $R_T$ is a P-position exactly when
\begin{equation}
  m\text{ is odd},\quad\text{or}\quad \Grundy(P_m)=0,\quad\text{or}\quad 2m\mid q\mp1 .
  \label{eq:pair-P}
\end{equation}
\end{corollary}

\begin{proof}
Immediate from Theorem~\ref{thm:pair-closed}.
\end{proof}

\subsection{Density among admissible prime fields}

For fixed $(G,T)$ the projective-line value is determined by $q\bmod 2|G|$ together with the bounded
list of exceptional orbit indicators. This follows from the orbit equation
\begin{equation}
  q+1=|G|f+\sum_K a_K[G:K],
  \label{eq:general-orbit-equation}
\end{equation}
which determines $f\bmod2$. For the dihedral family the exceptional indicators are controlled by
$4n\mid q\mp1$, so period $8n$ suffices.

\begin{theorem}[Density one half]
\label{thm:density}
Fix $n\ge2$, a legal triple type, and either the split or the nonsplit torus family. Among
admissible prime fields, the relative Dirichlet density of P-positions is $\tfrac12$, and the
relative density of N-positions is also $\tfrac12$.
\end{theorem}

\begin{proof}
Let $h=(q-1)/2n$ in the split case and $h=(q+1)/2n$ in the nonsplit case. The four possibilities
$h\bmod4$ give four admissible reduced residue classes modulo $8n$: the corresponding prime
candidate is $q=2nh\pm1$, which is odd, and any common divisor with $n$ would divide $\pm1$, so
$\gcd(q,8n)=1$ for all four values of $h\bmod4$. Corollary~\ref{cor:pn-congruences} selects exactly
two of these four classes in every triple type. The prime number theorem for arithmetic
progressions, equivalently its standard Dirichlet-density consequence \cite{Davenport2000}, gives
the same Dirichlet density to each reduced residue class. Relative to primes in the chosen torus
family the four classes therefore each have density $1/4$, proving the claim.
\end{proof}

\phasenote{C307/C308}{Theorem~\ref{thm:density} is one of the paper's two Lean-anchored headline
results and must keep that status (Fable's change~1). C307 attaches the C278 boundary sentence: the
density-$\tfrac12$ kernel certificate is conditional on exactly one quarantined equidistribution
axiom (\texttt{primes\_equidistribute}), audited as
\texttt{[propext, Classical.choice, Quot.sound, primes\_equidistribute]}. The paper must not imply
an axiom-free Lean proof. C307 also re-derives the statement against the corrected
Corollary~\ref{cor:pn-congruences}.}

\begin{corollary}[Density for the pair family]
\label{cor:pair-density}
Fix $m\ge3$ and a torus family. Among admissible prime fields the relative Dirichlet density of
P-positions for the two-point game is
\begin{equation}
  \begin{cases}
    1, & m\text{ odd, or } m\text{ even with }\Grundy(P_m)=0,\\[2pt]
    \tfrac12, & m\text{ even with }\Grundy(P_m)\ne0 .
  \end{cases}
  \label{eq:pair-density}
\end{equation}
\end{corollary}

\begin{proof}
An odd $m$ or a Dawson zero $\Grundy(P_m)=0$ makes every admissible field a P-position by
Corollary~\ref{cor:pair-P}. Otherwise, with $h=(q\mp1)/m$, the indicator
$\delta=\mathbf 1_{2m\mid q\mp1}$ equals $\mathbf 1_{2\mid h}$; the two parities of $h$ correspond
to two reduced residue classes modulo $2m$, each of equal Dirichlet density among primes, so
P-positions have relative density $\tfrac12$.
\end{proof}

\subsection{Worked examples}

\begin{example}[$D_{12}$]
\label{ex:d12}
Take $n=3$, so $G=D_{12}$ has rotation order six, and there are two legal triple types.

For odd $d$, each free orbit contributes $\Grundy(\Mob_{12})=1$ while every reflection orbit
contributes $\Grundy(L_2)=0$; the value is simply $f\bmod2$.

For even $d$, each free orbit is the prism $\Cyc_6\square K_2$ and contributes zero, while each
split reflection class contributes a ladder $L_3$ or $L_1$, both of value one; the total value is
$a_0\oplus a_1$. In the split-torus family $6\mid q-1$, the second reflection class is split exactly
when $12\mid q-1$, so the even-$d$ position is P exactly when $12\mid q-1$.
\end{example}

\phasenote{C307}{Fable's answer~(6): keep the $D_{12}$ example and add one contrasting polyhedral
example, drawn from the $A_5$ rows, rather than replacing it. That example depends on the C290
congruence laws and is therefore owed by C307, not by the structural rebuild.}

\section{Realization, boundaries, and outlook}
\label{sec:boundaries}

\subsection{Converse realization}

\begin{theorem}[Realization, three-point family]
\label{thm:realization-triple}
For the Klein-four boundary, $K_4$ occurs over infinitely many prime fields as a transitive residual
template of a legal tame conic-involution triple. For every $n\ge2$, the graphs
\begin{equation}
  \Mob_{4n},\qquad L_{n-1}
  \label{eq:realized-odd}
\end{equation}
and the completely deleted rotation-torus template occur over infinitely many prime fields as
transitive residual templates of legal tame dihedral conic-involution triples. If $n\ge3$ is odd,
the same is true of
\begin{equation}
  \Cyc_{2n}\square K_2,\qquad L_n,\qquad L_{n-2}.
  \label{eq:realized-even}
\end{equation}
\end{theorem}

\begin{proof}
The abstract coset calculations of Section~\ref{sec:triples} produce the listed graph whenever the
corresponding orbit type occurs.

For the Klein-four assertion, take primes $q\equiv3\pmod4$ and the standard commuting triple
\[
  t\mapsto-t,\qquad t\mapsto a/t,\qquad t\mapsto-a/t
\]
with $a$ a square. Its centres form a legal self-polar triangle. Exactly two of the three
involutions are split, so Theorem~\ref{thm:v4} gives $(q-3)/4$ free $K_4$-orbits, which is positive
for $q\ge7$. Dirichlet's theorem supplies infinitely many such primes.

Choose primes $q\equiv1\pmod{2n}$ for split-torus realizations or $q\equiv-1\pmod{2n}$ for
nonsplit-torus realizations; stronger congruences modulo $4n$ realize both reflection classes, and
Dirichlet's theorem supplies infinitely many such primes. The triples $T_d$ are legal by
Section~\ref{sec:triples}. For sufficiently large primes in any chosen class,
equation~\eqref{eq:orbit-equation} gives $f>0$, so the free template occurs. Reflection and
rotation-torus templates occur under the splitting conditions of Section~\ref{sec:arithmetic}.
\end{proof}

\begin{theorem}[Realization, two-point family]
\label{thm:realization-pair}
For every $m\ge3$, the cycle $\Cyc_{2m}$ and the path $P_m$ occur as transitive residual templates
of legal tame conic-involution pairs over infinitely many prime fields.
\end{theorem}

\begin{proof}
Choosing $q\equiv-1\pmod m$ (nonsplit) or $q\equiv1\pmod m$ (split) realizes a $D_{2m}$-action with
$m=\ord(\sigma_x\sigma_y)$. For large such primes the orbit count gives a free orbit, hence
$\Cyc_{2m}$, while a congruence modulo $2m$ that makes a reflection class split produces the
reflection orbit $P_m$. Dirichlet's theorem supplies infinitely many primes in each class.
\end{proof}

\subsection{The wild boundary}
\label{subsec:wild}

Every result above assumes the tame hypothesis $p\nmid 2m$. When $p\mid 2m$ --- equivalently, since
$q$ is odd, $p\mid m$ --- the classification changes character completely, for one structural
reason. A cyclic subgroup of $\PGL_2(q)$ whose order is divisible by $p$ has order exactly $p$,
because a nonsemisimple element of $\PGL_2(q)$ is unipotent of order $p$. Hence the rotation
$r=\sigma_x\sigma_y$ can never have order $2p,3p,\dots$, and the only wild dihedral group that
arises is $D_{2p}$, with $r$ a transvection.

Three consequences follow. First, $r$ has a single fixed point on $\PP^1(q)$ instead of the tame $0$
or $2$, so the deleted set $D_S=\Fix(r)$ is a single point and the even orbit-deletion pattern of
Section~\ref{sec:pairs} fails. Second, that fixed point is fixed by the whole group: $D_{2p}$ lies
in a Borel subgroup and acts reducibly on $\PP^1(q)$, so there is no regular orbit and no cycle
template $\Cyc_{2p}$; the action splits into the deleted fixed point and one orbit of size $p$ with
stabilizer of order $2$. The residual is therefore a single path,
\[
  R_S\cong P_p,\qquad \Grundy(R_S)=\Grundy(P_p)=\mathrm{A}002187(p).
\]
Third, because $m=p$ is odd, this directly contradicts the tame law that odd-order dihedral games are
P-positions (Theorem~\ref{thm:pair-closed}): the wild game $D_{2p}$ is generically an N-position. The
smallest odd nontrivial case $p=5$ gives $D_{10}$ with $r$ unipotent, $D_S$ a single point, residual
$P_5$, and $\Grundy=\mathrm{A}002187(5)=3\ne0$ --- an N-position, whereas every tame $D_{10}$ is a
P-position. Legal triples cannot be wild at all: a wild rotation is unipotent of odd order and has
no central involution, so no legal generating triple exists.

\phasenote{C307}{Fable's change~3: this wild result is the paper's sharpest boundary marker and
should be presented as such rather than as an aside. Whether to promote it from remark to a numbered
lemma is a user scope call; C283 judged it a short lemma and the default is the remark as landed.}

\subsection{Boundaries of the classification}

Three boundaries should be distinguished sharply.

\begin{enumerate}[label=(\arabic*)]
  \item \textbf{Completed here.} The tame dihedral configuration families of every rotation order
    and both configuration sizes, together with the finite polyhedral boundary of
    Section~\ref{sec:polyhedral}.
  \item \textbf{Wild polyhedral characteristic.} Subsection~\ref{subsec:wild} settles the wild
    dihedral pair; the wild polyhedral case is not treated here.
  \item \textbf{Growing full or subfield groups.} In the motivating projective-cap game, a generic
    fourth move typically generates the full group $\PSL_2(q)$ or $\PGL_2(q)$. Evaluating or
    controlling those growing residual templates is a separate problem, and no claim about it is
    made here.
\end{enumerate}

\subsection{Discussion}

The classification of legal dihedral configurations is complete for every rotation order and both
configuration sizes, and has a simple conceptual form. A legal tame dihedral conic configuration
produces only $K_4$'s, Möbius ladders, prisms, ordinary ladders, cycles, paths, and completely
deleted exceptional orbits. Its conic-only Node-Kayles value is an explicit congruence function of
the field and the embedded configuration. Three selected points give the even-rotation groups
$D_{4n}$ and values in $\{0,1\}$; two selected points give every $D_{2m}$, with Dawson's-chess
values, and the odd-order groups are always P-positions.

The fixed-point-deleted Schreier construction is what separates geometry from game theory. The
abstract pair $(G,T)$ determines a finite table of transitive template nimbers; the ambient geometry
determines only the parity vector of orbit multiplicities. The three-point and two-point families
share this reduction and differ only in that table: ladders and Möbius ladders for the former, paths
and cycles for the latter.

\phasenote{C311}{The Discussion is rewritten last, after the theorem boundary and the final title
are fixed (runbook migration table, current \S15 row).}

\appendix

\section{Evidence and validation map}
\label{app:evidence}

\phasenote{C307/C308}{This appendix is a structural placeholder. C307 populates it with the C281
dihedral census (all tame legal pairs and triples over prime $q\le23$) and the C288 polyhedral
embedding census (odd prime powers $q\le101$), together with the polyhedral edge lists and replay
counts displaced from Section~\ref{sec:polyhedral}. C308 adds the mixed-verification claim-route
map, the Lean adequacy statement, and the provenance boundary. C309 adds the regeneration commands
and manifest hashes. Censuses are validation evidence, not narrative spine.}

The path and cycle Node-Kayles values used in Section~\ref{sec:pairs}, the even-index property of
the Dawson zeros underlying Lemma~\ref{lem:cycle-zero}, and the full two-point classification
of~\eqref{eq:pair-templates}, \eqref{eq:pair-orbit} and~\eqref{eq:pair-closed} were verified end to
end by an exhaustive conic simulation over the prime fields $q\in\{5,7,11,13,17,19,23\}$: for every
legal pair the fixed-point-deleted residual graph was built directly and its Node-Kayles value
matched the orbit formula, with no discrepancy.

\section{Regular polyhedral templates: exact computation}
\label{app:regular-computation}

The computation behind Table~\ref{tab:polyhedral-regular} uses memoized Sprague-Grundy recursion
with connected-component xor and canonical masks under graph automorphisms. Left multiplication
supplies $|G|$ automorphisms in every class. The symmetric $(5,5,5)$ triple also admits the
colour-permuting $S_3$, giving a $360$-element canonicalization group.

The four $S_4$ values were cross-checked by two independent solvers. The five $A_5$ values were
subsequently independently cross-checked as well: a separate solver using a different
canonicalization reproduces all five, and an independent graph-isomorphism canonicalization confirms
that the automorphism group used equals the full automorphism group of each Cayley graph.

These are free-orbit values only. In particular the uniform value zero for regular $S_4$ orbits does
not evaluate a generic fourth-move residual whose generated group has grown to full $\PSL_2(q)$ or
$\PGL_2(q)$.

\begin{thebibliography}{99}

\bibitem{BrownEtAl2020}
S.~Brown, S.~Daugherty, E.~Fiorini, B.~Maldonado, D.~Manzano-Ruiz, S.~Rainville, R.~Waechter, and
T.~W.~H. Wong,
\newblock Nimber sequences of Node-Kayles games,
\newblock \emph{Journal of Integer Sequences} \textbf{23} (2020), Article 20.3.5.

\bibitem{Beauville2010}
A.~Beauville,
\newblock Finite subgroups of $\PGL_2(K)$,
\newblock in \emph{Vector Bundles and Complex Geometry}, Contemporary Mathematics \textbf{522}
(2010), 23--29.

\bibitem{HugganHuntemannStevens2021}
M.~A. Huggan, S.~Huntemann, and B.~Stevens,
\newblock The combinatorial game Nofil played on Steiner triple systems,
\newblock 2021, arXiv:2103.13501.

\bibitem{Tranchida2025}
P.~Tranchida,
\newblock Triples of involutions in $\PGL(2,q)$ and their incidence geometries,
\newblock \emph{Innovations in Incidence Geometry} \textbf{22} (2025).

\bibitem{Kobayashi2020}
Y.~Kobayashi, On structural parameterizations of Node Kayles, arXiv:2003.11775, 2020.

\bibitem{Dickson1901}
L.~E. Dickson,
\newblock \emph{Linear Groups with an Exposition of the Galois Field Theory},
\newblock Teubner, 1901.

\bibitem{Davenport2000}
H.~Davenport,
\newblock \emph{Multiplicative Number Theory}, 3rd ed., revised by H.~L. Montgomery,
\newblock Graduate Texts in Mathematics \textbf{74}, Springer, 2000.

\bibitem{WinningWays}
E.~R. Berlekamp, J.~H. Conway, and R.~K. Guy,
\newblock \emph{Winning Ways for Your Mathematical Plays}, 2nd ed.,
\newblock A K Peters, 2001--2004. (Dawson's chess, octal game $0.137$.)

\bibitem{OEIS-A002187}
OEIS Foundation,
\newblock A002187: Sprague-Grundy values for Dawson's chess,
\newblock \emph{The On-Line Encyclopedia of Integer Sequences},
\newblock \url{https://oeis.org/A002187}.

\end{thebibliography}

\end{document}
